\section*{Capítulo \ref{ch:g11:dissolution} --- La disolución de sólidos iónicos y moleculares}

\begin{solution}{exo:g11:dissolution:1}
\ce{NaCl(s) -> Na+(aq) + Cl-(aq)};
\ce{CaCl2(s) -> Ca^{2+}(aq) + 2Cl-(aq)};
\ce{Na2SO4(s) -> 2Na+(aq) + SO4^{2-}(aq)};
\ce{AlCl3(s) -> Al^{3+}(aq) + 3Cl-(aq)}.
\end{solution}

\begin{solution}{exo:g11:dissolution:2}
$[\ce{Na+}] = 2 \times 0.050 = \qty{0.10}{mol/L}$;
$[\ce{SO4^{2-}}] = \qty{0.050}{mol/L}$.
\end{solution}

\begin{solution}{exo:g11:dissolution:3}
La \omterm{def:g11:dissolution:solvation}{hidratación} es el hecho de que el \omterm{def:g9:ions:ion}{ion} quede rodeado de \omterm{def:g7:atoms-and-molecules:molecule}{moléculas} de
agua unidas a él. El extremo del oxígeno ($\delta^-$) apunta hacia un
\omterm{def:g9:ions:ion}{catión}, y el extremo de los hidrógenos ($\delta^+$), hacia un \omterm{def:g9:ions:ion}{anión}: las
cargas opuestas se atraen.
\end{solution}

\begin{solution}{exo:g11:dissolution:4}
$M(\ce{AlCl3}) = 27.0 + 3 \times 35.5 = \qty{133.5}{g/mol}$;
$n = 2.67 / 133.5 = \qty{0.0200}{mol}$; $c = 0.0200 / 0.2000 =
\qty{0.100}{mol/L}$; $[\ce{Al^{3+}}] = \qty{0.100}{mol/L}$,
$[\ce{Cl-}] = \qty{0.300}{mol/L}$.
\end{solution}

\begin{solution}{exo:g11:dissolution:5}
La cabeza \ce{-COO-} es iónica y atraída por el agua: \omterm{def:g11:dissolution:hydrophilic}{hidrófila}. La cola
de 17 \omterm{def:g7:atoms-and-molecules:atom}{átomos} de carbono, con solo enlaces \ce{C-C} y enlaces \ce{C-H}
casi apolares, es apolar: \omterm{def:g11:dissolution:hydrophilic}{hidrófoba}.
\end{solution}

\begin{solution}{exo:g11:dissolution:6}
La sal: agua (sólido iónico, \omterm{def:g6:solutions-and-solubility:solute}{disolvente} polar). La cera: ciclohexano
(cadenas apolares). El azúcar: agua (sus grupos \ce{O-H} forman \omterm{def:g11:polarity-and-cohesion:hydrogen-bond}{enlaces
de hidrógeno} con el agua). El yodo: ciclohexano (\omterm{def:g11:polarity-and-cohesion:polar-molecule}{molécula apolar}).
\end{solution}

\begin{solution}{exo:g11:dissolution:7}
El grupo \ce{O-H} del etanol forma \omterm{def:g11:polarity-and-cohesion:hydrogen-bond}{enlaces de hidrógeno} con las
\omterm{def:g7:atoms-and-molecules:molecule}{moléculas} de agua, que sustituyen a los que rompe: el etanol encaja en
el agua. El hexano, apolar, no puede formar enlaces así; las \omterm{def:g7:atoms-and-molecules:molecule}{moléculas}
de agua, unidas entre sí por \omterm{def:g11:polarity-and-cohesion:hydrogen-bond}{enlaces de hidrógeno}, no le dejan entrar.
\end{solution}

\begin{solution}{exo:g11:dissolution:8}
Las colas \omterm{def:g11:dissolution:hydrophilic}{hidrófobas} evitan el agua y se \omterm{def:g2:mixing-and-dissolving:dissolve}{disuelven} en la grasa, en el
centro; las cabezas cargadas son atraídas por el agua y se quedan fuera.
Todas las \omterm{def:g11:dissolution:amphiphilic}{micelas} llevan cargas negativas en su superficie, y las cargas
del mismo signo se repelen: se mantienen separadas.
\end{solution}

\begin{solution}{exo:g11:dissolution:9}
El ciclohexano: es \omterm{def:g6:solutions-and-solubility:miscible}{inmiscible} con el agua, así que forma una capa aparte
que se puede separar, y el aroma (apolar) se \omterm{def:g2:mixing-and-dissolving:dissolve}{disuelve} bien en él. El
etanol es \omterm{def:g6:solutions-and-solubility:miscible}{miscible} con el agua: no se forma una segunda capa y no se
puede separar nada.
\end{solution}

\begin{solution}{exo:g11:dissolution:10}
El ciclohexano es menos denso que el agua. Se vacía por la llave la capa
inferior, la acuosa. El yodo está en la capa superior de ciclohexano,
que se queda en el embudo.
\end{solution}

\begin{solution}{exo:g11:dissolution:11}
Nada: el sulfato de cobre es iónico y no se \omterm{def:g2:mixing-and-dissolving:dissolve}{disuelve} en el ciclohexano,
apolar. El agua sigue azul y el ciclohexano, incoloro.
\end{solution}

\begin{solution}{exo:g11:dissolution:12}
$[\ce{Cl-}] = 2c$, así que $c = \qty{0.150}{mol/L}$;
$n = 0.150 \times 0.2500 = \qty{0.0375}{mol}$;
$m = 0.0375 \times 111.1 = \qty{4.17}{g}$.
\end{solution}

\begin{solution}{exo:g11:dissolution:13}
Los \omterm{def:g9:ions:ion}{iones} calcio y magnesio precipitan los \omterm{def:g9:ions:ion}{iones} estearato como jabón de
cal: el jabón se gasta antes de poder formar \omterm{def:g11:dissolution:amphiphilic}{micelas}, y hace poca
espuma. Solo el calcio: \qty{0.010}{mol} de \ce{Ca^{2+}} se llevan
\qty{0.020}{mol} de estearato, es decir, $0.020 \times 306.0 =
\qty{6.1}{g}$ de estearato de sodio por litro.
\end{solution}

\begin{solution}{exo:g11:dissolution:14}
Volumen: \qty{0.2000}{L}. \ce{Na+}: $0.020 / 0.2000 = \qty{0.10}{mol/L}$;
\ce{Ca^{2+}}: $0.010 / 0.2000 = \qty{0.050}{mol/L}$; \ce{Cl-}:
$(0.020 + 0.020) / 0.2000 = \qty{0.20}{mol/L}$. Carga positiva:
$0.10 + 2 \times 0.050 = 0.20$; negativa: $0.20$. Neutra.
\end{solution}

\begin{solution}{exo:g11:dissolution:15}
Un litro de heptano tiene una masa de \qty{0.68}{kg} y puede contener
$17 \times 0.68 = \qty{12}{g}$ de yodo, unas 40 veces más que un litro
de agua. Al agitarlos juntos, casi todo el yodo pasa al heptano: la
\omterm{def:g10:chemical-species:extraction}{extracción} es eficaz.
\end{solution}

\begin{solution}{pb:g11:dissolution:1}
\textbf{1.} De las rocas por las que ha pasado el agua (caliza,
dolomía, yeso), que se \omterm{def:g2:mixing-and-dissolving:dissolve}{disuelven} un poco.

\textbf{2.} $[\ce{Ca^{2+}}] = 0.120 / 40.1 = \qty{2.99e-3}{mol/L}$;
$[\ce{Mg^{2+}}] = 0.024 / 24.3 = \qty{9.9e-4}{mol/L}$.

\textbf{3.} \ce{Ca(HCO3)2 -> Ca^{2+}(aq) + 2HCO3-(aq)}; así que
$[\ce{HCO3-}] = 2 \times \num{2.99e-3} = \qty{5.99e-3}{mol/L}$.

\textbf{4.} $\num{2.99e-3} \times 150 = \qty{0.449}{mol}$.

\textbf{5.} $18 \times 12.0 + 35 \times 1.0 + 2 \times 16.0 + 23.0 =
\qty{306.0}{g/mol}$.

\textbf{6.} \ce{C17H35COONa(s) -> C17H35COO-(aq) + Na+(aq)}.

\textbf{7.} Cabeza \ce{-COO-}: \omterm{def:g11:dissolution:hydrophilic}{hidrófila}; cola \ce{C17H35-}: \omterm{def:g11:dissolution:hydrophilic}{hidrófoba}.
Tiene los dos tipos de parte: es \omterm{def:g11:dissolution:amphiphilic}{anfífilo}.

\textbf{8.} Una esfera diminuta de \omterm{def:g9:ions:ion}{iones} de jabón, con las colas hacia
dentro y las cabezas hacia fuera. Las colas se \omterm{def:g2:mixing-and-dissolving:dissolve}{disuelven} en la grasa,
que se rompe en gotitas envueltas en \omterm{def:g11:dissolution:amphiphilic}{micelas}; las superficies cargadas
mantienen separadas las gotitas en el agua, que se las lleva en el
aclarado.

\textbf{9.} El jabón tiene que formar \omterm{def:g9:ions:ion}{iones} y \omterm{def:g11:dissolution:amphiphilic}{micelas}, lo que necesita
agua alrededor de las cabezas; y el agua del aclarado es la que se lleva
la grasa. En aceite, la grasa se quedaría simplemente disuelta sobre la
piel.

\textbf{10.} \ce{2C17H35COO- + Ca^{2+} -> Ca(C17H35COO)2}: cargas
$2 \times (-1) + 2 = 0$ a la izquierda, y $0$ a la derecha.

\textbf{11.} \ce{C17H35COO-}: exceso $- 2x$; \ce{Ca^{2+}}:
$\num{2.99e-3} - x$; estearato de calcio: $x$. Los \omterm{def:g9:ions:ion}{iones} calcio son el
limitante: $x_{\max} = \qty{2.99e-3}{mol}$.

\textbf{12.} $2 \times \num{2.99e-3} = \qty{5.99e-3}{mol}$ por litro.

\textbf{13.} $M = 36 \times 12.0 + 70 \times 1.0 + 4 \times 16.0 + 40.1 =
\qty{606.1}{g/mol}$; $\num{2.99e-3} \times 606.1 = \qty{1.81}{g}$ de
costra por litro.

\textbf{14.} $2 \times \num{9.9e-4} = \qty{1.98e-3}{mol}$ por litro.

\textbf{15.} $2 \times 0.449 = \qty{0.898}{mol}$.

\textbf{16.} $0.898 \times 306.0 = \qty{275}{g}$.

\textbf{17.} Magnesio: $\num{9.9e-4} \times 150 = \qty{0.148}{mol}$, que
se llevan \qty{0.296}{mol} de estearato, \qty{91}{g} de jabón; en total,
unos $275 + 91 = \qty{366}{g}$.

\textbf{18.} $275 / 100 \approx 2.7$ pastillas.

\textbf{19.} El calcio de un baño desperdicia unos \qty{0.27}{kg} de
jabón en forma de costra.
\end{solution}
